Taken together, the electrode-level comparisons in Figure~\ref{fig:exp1_subset_per_electrode_stats} show that a visibly lower descriptive mean did not by itself imply a statistically significant change. Fp1 and Fp2, the highest-performing electrode in each dataset, showed no detectable paired difference after correction in either corpus, while most other electrodes changed significantly, generally by a few percentage points. Significant differences were also more widespread on Cao2018 than on Internal, where roughly half of the non-frontal electrodes did not reach significance despite showing the same direction of change.
%These tests probe a configuration effect, whether an electrode's own detection performance changes when its region operates alone rather than as part of the complete montage. That question is distinct from the spatial pattern described above, in which frontal electrodes simply detect blinks better than posterior electrodes regardless of which subset scored them.
